% !TeX root = ../example-en.tex
\section{Callouts}
\sectioncoverpage[\campusexamplesectionthreeimage]{Callout styles, semantic roles and usage.}
\note{Use a small amount of visual emphasis to distinguish advice, context and warnings.}
\subsection{Styles}
\begin{frame}{Tips, notes and alerts}
  \begin{tipbox}[Tip]
    Use \texttt{tipbox} for practical advice, such as a recommended experiment setup.
  \end{tipbox}
  \begin{notebox}[Note]
    Use \texttt{notebox} for context, such as data sources or experimental conditions.
  \end{notebox}
  \begin{alertbox}[Alert]
    Use \texttt{alertbox} for pitfalls, such as settings that matter for reproduction.
  \end{alertbox}
\end{frame}
\note{Compare the functional colors. Each box should carry one point that deserves emphasis.}

\begin{frame}{Content and production}
  \begin{columns}[T,onlytextwidth]
    \begin{column}{.46\textwidth}
      \textbf{Outline stage}
      \begin{itemize}
        \item Organize sections and topics
        \item Develop explanations and evidence
        \item Review and approve the content
      \end{itemize}
    \end{column}
    \begin{column}{.46\textwidth}
      \textbf{Production stage}
      \begin{itemize}
        \item Divide slides and choose layouts
        \item Allocate body text and notes
        \item Validate and export PDF/PPTX
      \end{itemize}
    \end{column}
  \end{columns}
  \begin{notebox}[Before production]
    Approve the detailed outline; keep essential evidence and limitations visible.
  \end{notebox}
\end{frame}
\note{Compare the two stages before discussing the condition below them. A note box suits necessary context, a tip box suits advice, and an alert box suits a critical limitation. Do not put every sentence in a box or repeat the slide's entire message.}

\subsection{Definitions and examples}
\begin{frame}{Definitions and examples}
  \begin{definitionbox}[Definition]
    A \textbf{modality} is a form of information, such as text, images or audio.
    Multimodal learning models these heterogeneous sources jointly.
  \end{definitionbox}
  \begin{examplebox}[Example]
    Tensor Fusion uses an outer product to retain unimodal, bimodal and trimodal interactions explicitly, forming a richer joint representation.
  \end{examplebox}
\end{frame}
\note{Distinguish definitions from examples. The Tensor Fusion description is retained from the Chinese template example; this slide demonstrates semantic boxes.}

\subsection{Usage}
\begin{frame}[fragile]{Callout syntax}
  \begin{verbatim}
\begin{tipbox}[Custom title]
  The point to emphasize.
\end{tipbox}
  \end{verbatim}
  \begin{itemize}
    \item Other environments: \texttt{notebox}, \texttt{alertbox}, \texttt{examplebox}, \texttt{definitionbox}.
    \item Omit a title to use the default label; use \texttt{campusbox} for no title.
  \end{itemize}
  \begin{tipbox}
    Omitting the optional title uses the default label.
  \end{tipbox}
\end{frame}
\note{Default labels follow the document's language setting. The optional argument overrides an individual box title. For campusbox it selects the color instead, with tipcolor as the default; no title area is reserved.}


\begin{frame}[fragile]{Titleless content boxes}
  \begin{verbatim}
\begin{campusbox}
  Emphasize one point at a time.
\end{campusbox}
  \end{verbatim}
  \begin{campusbox}
    Emphasize one point at a time.
  \end{campusbox}
  \begin{campusbox}[notecolor]
    Select a color with the optional argument, such as \texttt{[notecolor]}.
  \end{campusbox}
  \begin{campusbox}[alertcolor]
    The title area is removed; the body keeps its colored left border.
  \end{campusbox}
\end{frame}
\note{All three boxes are titleless. campusbox uses tipcolor by default and accepts a color, not a title, as its optional argument. The code produces the first box; the next two demonstrate note and alert colors.}
